\documentclass[a4paper]{article}
\usepackage{graphicx} % Required for inserting images
\usepackage{amsmath}
\usepackage{amssymb}
\usepackage{amsthm}
\usepackage[most]{tcolorbox}
\usepackage[margin=0.75in]{geometry}

\title{Ranks of Cohen-Macaulay Modules}
\author{Jaspreet Singh Sidhu}
\date{October 2026}

\begin{document}


\maketitle

Given a polynomial ring $R=k[x_{0},x_{1},x_{2},\dots,x_{n}]$ and a homogeneous polynomial $f\in R$, presence of non-free maximal Cohen-Macaulay modules (MCMs) over $R/(f)$ indicates existence of a singularity, by Auslander-Buchsbaum formula. Note that the affine cone of $f$ has a singularity at the vertex $\mathfrak{m}$. Let $Q=\widehat{P_{\mathfrak{m}}}$, $S=Q/(f)$, and $e=\lfloor n/2\rfloor$. As per Buchweitz-Greuel-Schreyer's (BGS) conjecture, in \cite{BGS}, the rank of non-free MCMs over $S$ is bounded from below by $e$. In our work, we will prove the BGS conjecture in the case when $f$ has even degree by showing that the rank is positive multiple of constant depending upon $n$. This will in turn imply a weaker conjecture by Brown and Walker (see conjecture 1.3 in \cite{BrownWalker}) that given any nontrivial matrix factorization of $f$ consisting of finite-rank free $R$-modules $M$ and $M'$ together with maps $d:M\rightarrow M'$ and $d':M'\rightarrow M$ such that $dd'=f\cdot 1_{M'}$ and $d'd=f\cdot 1_{M}$, the rank of module $M$ is at least $2^{\lfloor \frac{n}{2}\rfloor}$. The connection, suggested by the two conjectures, is made clear by the equivalence (as triangulated categories) of the homotopy category of $\mathbb{Z}/{2}\mathbb{Z}$-graded matrix factorizations $\text{HMF}_{gr}(f)$ and the stable category of $\mathbb{Z}/{2}\mathbb{Z}$ MCMs over $R/(f)$. As of now, the conjectures have not been proven independently.\par
In addition to the above conjectures, there is also a graded problem asking for the lower bound on the rank of Ulrich bundles on the projective variety determined by $f$. An Ulrich bundle of rank $r$ on the projective variety $V(f)$ is a $R/(f)$ module $M$ with a linear $R$-free presentation:
\[
 0\longrightarrow R(-1)^{rd}\xrightarrow{\ \phi\ }R^{rd}
 \longrightarrow M\longrightarrow0;
\]
Earlier, in \cite{Dyckerhoff}, Dyckerhoff showed that starting with a local regular ring $R$ with residue field $k$ and maximal ideal $\mathfrak{m}$, given an element $f\in R$, the category $\operatorname{MF}^{\infty}(f)$ of matrix factorizations of $f$ of potentially infinite rank is compactly generated by a stabilization $k^{stab}$ of the residue field. The result of our interest is the quasiequivalence of $\operatorname{MF}^{\infty}(f)$ with an explicit model of the derived category of $\operatorname{End}(k^{stab})^{op}$-modules, where $\operatorname{End}(k^{\text{stab}})$ is computed in $\operatorname{MF}^{\infty}(f)$. This is a new setting in which to investigate the BGS and Brown-Walker's conjecture. Our objective is to translate and investigate the conjectures in this setting.\par

A more modest objective is that there is an $A_{\infty}$ algebra structure on the algebra $A=\operatorname{End}(k^{stab})^{op}$, which as explained in \cite{Dyckerhoff}, allows one to see that $MF^{\infty}(f)$ is quasiequivalent to the derived category of modules over $H^{*}(A)$. However, the $A_{\infty}$ structure is tedious to compute in practice, as in \cite{Dyckerhoff}. We will try to extend the known computations from the case of quadratic forms in $k[x_{1},x_{2},\dots,x_{n}]$ with $\operatorname{char}(k)\ne 2$ to hypersurfaces in general.\par


\newpage

\begin{tcolorbox}[
  colback=yellow!12,
  colframe=yellow!50!black,
  boxrule=0.6pt,
  arc=1mm,
  left=2mm,
  right=2mm,
  top=1.5mm,
  bottom=1.5mm,
  boxsep=0pt,
  before skip=0pt,
  after skip=0pt
]

Given a polynomial ring $R=k[x_{0},x_{1},x_{2},\dots,x_{n}]$ and a homogeneous polynomial $f\in R$, with an isolated singularity at the origin. 
The category of maximal Cohen modules over $R/(f)$ measure the complexity of the singularity. The minimal rank of non-trivial MCM in this category is the subject 
of a long standing conjectures, Buchweitz-Greuel-Schreyer's (BGS) conjecture, in \cite{BGS}. More recently, Brown and Walker have stated a variation that can 
be phrased at the level of derived categories. \textbf{provide a few more details of these conjectures here}

The goal of this work is study these conjectures and rephrase them in terms of $A_\infty$-categories. The latter question is highly tractible. 
The results of \cite{Dyckerhoff} provide a canonical generator $k^{\rm stab}$ for the stable ${\mathbb Z}/2$-graded category of MCMs. Derived Morita theory then furnishes an equivalence 
with the category of $A_\infty$-modules over the endomorphism ring of this generator. Hence the aforementioned conjectures have $A_\infty$ restatements that I will write down in my 
thesis.


\textbf{Note: }I would not stated that you are going to prove BGS without providing a concrete way to actually tackle the conjecture. An expert will know that you won't be able to complete this.

\end{tcolorbox}

\clearpage

\bibliography{bib.bib}
\bibliographystyle{plain}
\end{document}
